% !TEX root = IE3_expA_prelab.tex
\documentclass[lualatex,a4paper,10pt]{jlreq}
% Shared LuaLaTeX preamble.
\usepackage{luatexja-fontspec}
\setmainfont{TeX Gyre Termes}
\setsansfont{TeX Gyre Heros}
\setmonofont{Latin Modern Mono}
\setmainjfont[BoldFont=HaranoAjiGothic-Medium]{HaranoAjiMincho-Regular}
\setsansjfont[BoldFont=HaranoAjiGothic-Bold]{HaranoAjiGothic-Medium}
\usepackage[top=22mm,bottom=22mm,left=23mm,right=23mm]{geometry}
\usepackage{amsmath,amssymb,booktabs,array,longtable,tabularx}
\usepackage{graphicx,xcolor,tikz,circuitikz}
\usetikzlibrary{arrows.meta,positioning,calc}
\usepackage{enumitem,fancyhdr,listings,needspace}
\usepackage[unicode,colorlinks=true,linkcolor=labblue,urlcolor=labteal]{hyperref}
\usepackage{bookmark}
\definecolor{labblue}{RGB}{30,70,112}
\definecolor{labteal}{RGB}{0,100,105}
\definecolor{labgray}{gray}{0.95}
\ModifyHeading{section}{font={\large\sffamily\bfseries\color{labblue}}}
\ModifyHeading{subsection}{font={\normalsize\sffamily\bfseries\color{labteal}}}
\setlength{\parindent}{1\zw}
\setlength{\parskip}{3pt}
\setlength{\emergencystretch}{2em}
\renewcommand{\arraystretch}{1.35}
\setlist{itemsep=3pt,topsep=4pt,leftmargin=2\zw}
\newcolumntype{Y}{>{\raggedright\arraybackslash}X}
\newcolumntype{C}[1]{>{\centering\arraybackslash}p{#1}}
\pagestyle{fancy}
\fancyhf{}
\fancyhead[L]{\small\color{labblue} IE3 後期・電子工学実験}
\fancyhead[R]{\small テーマ\LabCode}
\fancyfoot[C]{\thepage}
\setlength{\headheight}{16pt}
\DeclareOldFontCommand{\rm}{\normalfont\rmfamily}{\mathrm}
\lstset{basicstyle=\small\ttfamily,columns=fullflexible,keepspaces=true,
 breaklines=true,frame=single,backgroundcolor=\color{labgray},
 numbers=left,numberstyle=\tiny,showstringspaces=false,aboveskip=8pt,belowskip=8pt}
\newcommand{\blank}{\rule{22mm}{0.3pt}}
\newcommand{\answerline}{\par\noindent\rule{\linewidth}{0.25pt}\par}
\newcommand{\writing}[1]{\par\Needspace{\dimexpr#1+5mm\relax}\noindent%
 \fcolorbox{labblue!35}{white}{\parbox[c][#1][t]{\dimexpr\linewidth-2\fboxsep-2\fboxrule\relax}{\vspace{2mm}\noindent\strut\hfill\par}}\par}
\newcommand{\hint}[1]{\par\smallskip\noindent{\sffamily\bfseries\color{labteal}記述の視点：}#1\par}
\newcommand{\note}[1]{\par\smallskip\noindent\fcolorbox{labblue!45}{labblue!5}{%
 \parbox{\dimexpr\linewidth-2\fboxsep-2\fboxrule\relax}{#1}}\par\smallskip}
\newcommand{\eqerror}{\varepsilon=\frac{x_{\mathrm{meas}}-x_{\mathrm{th}}}{x_{\mathrm{th}}}\times100\ \%}
\newcommand{\LabSource}{徳山工業高等専門学校 情報電子工学科，
 『2026年度 電子工学実験（3年後期）テキスト』，テーマ\LabCode，
 \expandafter\path\expandafter{\SourcePdf}。}
\newcommand{\reporttitle}{
 \noindent{\small\sffamily\color{labteal}実験報告書・記入ガイド}\par
 \noindent{\LARGE\sffamily\bfseries \LabCode\quad\LabTitle}\par\smallskip
 \noindent 氏名：\blank\quad 班：\blank\quad 実験日：\blank\par
 \note{目的・原理・手順を確認し、実測値と自分の考察を記入する。
 考察のガイドは、検討する事象・使う測定結果・記述の方針を示す。}
}
\newcommand{\prelabtitle}{
 \noindent{\small\sffamily\color{labteal}事前学習・前提知識プリント}\par
 \noindent{\LARGE\sffamily\bfseries \LabCode\quad\LabTitle}\par\smallskip
 \noindent 氏名：\blank\quad 班：\blank\quad 予習日：\blank\par
 \note{用語、回路の動き、計算、測定表への記入の順に読む。
 例題の値は説明用であり、実験結果ではない。}
}
\newcommand{\tocrow}[3]{\hyperref[#1]{#2}& #3 &\pageref*{#1}\\[2mm]}
\newcommand{\documentmap}[1]{
 \section*{目次と概要}
 \noindent 青色の項目をクリックすると該当箇所へ移動する。\par
 {\small\begin{longtable}{@{}p{0.29\linewidth}p{0.61\linewidth}r@{}}
 \toprule {\color{labblue}\bfseries 項目} & {\color{labblue}\bfseries 読むと分かること} & 頁\\\midrule
 \endhead #1\bottomrule
 \end{longtable}}
 \clearpage
}
\newenvironment{terms}{\par\smallskip\begin{longtable}{@{}p{0.23\linewidth}p{0.73\linewidth}@{}}
 \toprule {\color{labteal}\bfseries 用語}&{\color{labteal}\bfseries 意味・回路での役割}\\\midrule\endhead}
 {\bottomrule\end{longtable}}
\newcommand{\term}[2]{\textbf{#1}&#2\\[2mm]}
\newcommand{\guidepoint}[4]{
 \Needspace{32mm}\subsection{#1}
 \noindent{\bfseries\color{labteal}考える事象：}#2\par
 \noindent{\bfseries\color{labteal}参照する結果：}#3\par
 \noindent{\bfseries\color{labteal}記述の方針：}#4\par
}
\newcommand{\reportclosing}[1]{
 \clearpage\section{結論のまとめ方}\label{sec:closing}
 \hint{#1}
 \ConclusionGuide
 \begin{enumerate}
 \item 目的に対応する最も重要な実測結果を、測定条件と数値を添えて述べる。
 \item 原理と一致した点・一致しなかった点を示す。原因は確認済みの事実と推測を分ける。
 \item 実施範囲で言えることと、未確認の条件を一文ずつまとめる。
 \end{enumerate}
 \note{書き出しの型：「○○の条件で△△を測定したところ、□□となった。
 これは◇◇と整合する。ただし××は確認しておらず、一般化には追加測定が必要である。」
 空欄は自分の実測値と根拠で埋める。}
 \noindent 結論（実験後に記入）：\writing{30mm}
 \noindent 改善案・次に確認したい条件：\writing{16mm}
 \section*{参考文献}\label{sec:references}
 \begin{enumerate}
 \item \LabSource
 \item 使用したデータシート・書籍・ウェブ資料（著者、題名、該当箇所、URL、参照日）を追加する。
 \end{enumerate}
}
\newcommand{\prelabclosing}{
 \section*{準備・出典}\label{sec:prelabsource}
 \begin{itemize}[nosep]
 \item 資料、筆記具、記録用紙、電卓と、実験条件・機器設定の記録欄。
 \item 確認問題の解答と、分からない用語・回路点への具体的な質問。
 \end{itemize}
 {\small\noindent\LabSource\par}
}
\newcommand{\simpleflow}[1]{
 \begin{center}
 \begin{tikzpicture}[node distance=5mm and 5mm,
 every node/.style={font=\small},box/.style={draw=labblue,fill=labblue!4,rounded corners=1mm,
 align=center,minimum height=10mm,text width=30mm}]
 #1
 \end{tikzpicture}
 \end{center}
}

\newcommand{\LabCode}{A}
\newcommand{\LabTitle}{回路網}
\newcommand{\SourcePdf}{IE3_expA_A4.pdf}
\begin{document}
\prelabtitle
\documentmap{
\tocrow{sec:auto1}{1 用語を一つずつ理解する}{言葉の意味、単位、回路や測定での役割を確認する。}
\tocrow{sec:auto2}{2 端子から回路を読み直す}{部品と信号の経路を追い、状態が変化する理由を理解する。}
\tocrow{sec:auto3}{3 到達目標と基礎}{端子から見た等価回路と回路網定理の仕組みと成立条件を整理する。}
\tocrow{sec:auto4}{4 テブナン等価回路}{開放電圧と等価抵抗から負荷電流を予測する。}
\tocrow{sec:auto5}{5 例題：計器の影響}{電流計の内部抵抗による電流変化を計算する。}
\tocrow{sec:auto6}{6 3つの定理の使い分け}{重ねの理・相反・補償の適用条件を区別する。}
\tocrow{sec:auto7}{7 測定からレポートへ：順番に計算する}{実測値から計算・作図・記述へ進む順序を例題で練習する。}
\tocrow{sec:auto8}{8 確認問題}{自分で計算・説明して、実験に必要な理解を確かめる。}
\tocrow{sec:auto9}{解答の要点}{数値と理由を照合し、単位や論理の取り違えを見直す。}
\tocrow{sec:auto10}{実験の読み取り}{この項目の仕組みと実験で確認すべき点を整理する。}
\tocrow{sec:prelabsource}{準備・出典}{持参物と資料を確認し、具体的な質問を準備する。}
}

\section{用語を一つずつ理解する}\label{sec:auto1}
\begin{terms}
\term{回路網・端子}{抵抗や電源を接続した全体が回路網である。外部へつなぐ接続点が端子。内部が複雑でも、二つの端子から見た電圧と電流の関係を調べられる。}
\term{電位差と基準方向}{電圧は二点間の電位差、電流は選んだ矢印に沿う量である。負の測定値は「故障」ではなく、実際の向きが基準と逆であることを表す。}
\term{開放・短絡}{開放は枝をつながず電流が流れない状態。短絡は抵抗がほぼ零の導線でつなぐ状態。開放では電圧が零とは限らず、短絡では電流が零とは限らない。}
\term{等価回路}{指定した端子で同じ電圧・電流関係を示す別の回路。内部のすべての電圧や消費電力まで同じになるという意味ではない。}
\term{線形性}{入力を何倍かにすると応答も同じ倍率になり、入力を足すと応答も足せる性質。一定抵抗の回路が典型であり、重ねの理の前提となる。}
\term{計器の負荷効果}{電流計にも小さな直列抵抗、電圧計にも有限の入力抵抗がある。接続によって元の回路が変わるため、読みが理想値からずれることがある。}
\term{残差・相対誤差}{残差は「実測値−予測値」という差。相対誤差はその差を予測値で割った割合。予測値が零に近い場合には、割合だけでは比較しにくい。}
\end{terms}
\section{端子から回路を読み直す}\label{sec:auto2}
\subsection{テブナン等価を作る三段階}
まず負荷Bを外す。このとき端子の電流はほぼ零なので、測った端子電圧を開放電圧$V_{\rm th}$とする。
次に電源を外して独立電圧源の接続点を短絡し、端子から見た抵抗$R_{\rm th}$を測る。電流源がある理論問題では、独立電流源の作用を零にすると開放する。従属電源は無条件に消さない。
最後に負荷Bを戻す。電圧源から流れる電流は$R_{\rm th}$、$R_B$、電流計の$r_A$を順に通るので、分母には三つの抵抗を加える。
\note{電源を外した後に回路側の接続点を短絡する。動作中の電源出力そのものを短絡する操作とは異なる。抵抗を測るときも回路へ外部電圧を加えない。}
\subsection{重ねの理では向きをそろえる}
二電源を同時に接続した結果を、一電源ずつ作用させた結果の和で予測する。
二つの測定で電流計を逆につなぎ直すと符号が比較できない。回路図に矢印と電圧の正側を先に書いて、同じ基準の$I_1,I_2,V_1,V_2$を保存する。
電力は$I^2R$であり、電流を足してから二乗するため、各電源の電力をそのまま足すことはできない。
\subsection{相反定理では条件も交換実験の一部}
線形で双方向性のある抵抗回路で、電源を入れる枝と電流を観測する枝を交換する。
同じ電源電圧、同じ計器レンジ・内部抵抗で比較しないと、交換だけによる違いなのか判断しにくい。


\section{到達目標と基礎}\label{sec:auto3}
端子から見た回路を等価回路に直し、電源を1つずつ作用させて電流を求める。
電流計は直列、電圧計は並列につなぐ。抵抗測定は電源を外してから行う。
オームの法則 $V=RI$、節点の電流則、閉回路の電圧則を使う。
\[
R_s=R_1+R_2,\qquad R_p=\frac{R_1R_2}{R_1+R_2}.
\]
\section{テブナン等価回路}\label{sec:auto4}
負荷を外して開放端子電圧 $V_{\rm th}$ を求める。
独立電源をゼロにして、端子から見た $R_{\rm th}$ を求める。
負荷をつないだときの電流は $I=V_{\rm th}/(R_{\rm th}+R_L)$ になる。
\begin{figure}[h]\centering
\begin{circuitikz}
\draw (0,0) to[V,l=$V_{\rm th}$] (0,2.5) to[R,l=$R_{\rm th}$] (3,2.5)
to[R,l=$R_L$] (3,0) -- (0,0);
\end{circuitikz}\caption{テブナン等価回路（原理説明用）}\end{figure}
\section{例題：計器の影響}\label{sec:auto5}
$V_{\rm th}=6.0$ V、$R_{\rm th}=200\ \Omega$、$R_L=800\ \Omega$、
内部抵抗 $r_A=50\ \Omega$ とする。
\[
I=\frac{6.0}{200+800+50}=5.714\ {\rm mA}.
\]
計器を無視した6.000 mAより約4.76\%小さい。回路が変化した効果を先に見積もる。
\clearpage
\section{3つの定理の使い分け}\label{sec:auto6}
\begin{itemize}
\item 重ねの理：同じ正方向で電圧・電流を加える。$P=I^2R$ は直接加算できない。
\item 相反定理：線形かつ双方向性をもつ回路で励振と応答の位置を交換できる。増幅器には無条件で適用できない。
\item 補償定理：元の枝電流 $I$、抵抗増分 $\Delta R$ に対して $-I\Delta R$ の起電力を抵抗変更後の回路へ加え、電流変化を求める。
\end{itemize}
\section{測定からレポートへ：順番に計算する}\label{sec:auto7}
\subsection{テブナンの測定表}
\begin{enumerate}
\item 各電源条件の開放電圧と、電源を外して測った抵抗を別々に記録する。
\item すべての抵抗を$\Omega$、電圧をVにそろえ、予測電流をAで計算する。
\item 表の単位がmAなら1000倍して記入し、実測値との差を求める。
\end{enumerate}
\textbf{説明用の例：}$V_{\rm th}=4.0\,\mathrm{V}$、$R_{\rm th}=200\,\Omega$、$R_B=780\,\Omega$、$r_A=20\,\Omega$と仮定する。
\[
 I_{\rm pred}=\frac{4.0}{200+780+20}=0.0040\,\mathrm{A}=4.0\,\mathrm{mA}.
\]
計器を無視すると$4.0/980=4.082\,\mathrm{mA}$となる。
この二つを比べれば、約$0.082\,\mathrm{mA}$の差が計器の抵抗だけでも生じ得ると分かる。
実験では仮の数値を転記せず、使用レンジの内部抵抗で計算する。
\subsection{重ねの理の残差}
仮に$I_1=2.4\,\mathrm{mA},I_2=-1.0\,\mathrm{mA}$なら予測は$1.4\,\mathrm{mA}$。
同時接続の仮の測定値が$1.36\,\mathrm{mA}$なら、
$\Delta I=1.36-1.40=-0.04\,\mathrm{mA}$。
「少し違う」で終えず、電流計の最小目盛、読み取り幅、電源の再設定による変動とこの差を比較する。
\subsection{考察で使うメモ}
同じ電源条件を表の同じ行に並べ、抵抗・計器レンジ・測定方向を表の近くに示す。
比較できた条件を明記してから「理論と整合する」と書く。数行の測定だけで、あらゆる回路で成立したと結論づけない。

\clearpage
\clearpage
\section{確認問題}\label{sec:auto8}
\begin{enumerate}
\item $V_{\rm th}=5$ V、$R_{\rm th}=100\ \Omega$、負荷400 $\Omega$ の電流は？\answerline
\item 理想電圧源／理想電流源をゼロにした回路をそれぞれ描け。\answerline
\item 同じ基準方向で $I_1=3$ mA、$I_2=-1$ mAなら両電源時の電流は？\answerline
\item 抵抗の変更後の電流と、補償定理で求める電流変化は同じか。\answerline
\end{enumerate}
\section*{解答の要点}\label{sec:auto9}
1：10 mA。2：電圧源は短絡、電流源は開放。3：2 mA。
4：変更後の電流は元の電流に変化分を加えた値である。
\section*{実験の読み取り}\label{sec:auto10}
資料A-4〜A-6の実験手順を読み、端子名と正方向を回路図へ書き込む。
電流計の内部抵抗は資料の例示値だけで決めず、当日の計器・レンジを確認する。
\prelabclosing

\end{document}
